\documentclass[11pt]{article}

\usepackage[review]{acl}

\usepackage{times}
\usepackage{latexsym}
\usepackage[T1]{fontenc}
\usepackage[utf8]{inputenc}
\usepackage{microtype}
\usepackage{graphicx}
\usepackage{booktabs}
\usepackage{amsmath,amssymb}
\usepackage{xcolor}
\usepackage{multirow}
\usepackage{array}

\title{The Likert Blindspot: Detecting Hidden Style Bias\\in LLM Judges via Controlled Style Decomposition}

\author{Anonymous}

\begin{document}
\maketitle

\begin{abstract}
LLM-as-judge evaluation is increasingly central to NLP benchmarking and alignment, yet style biases---systematic preferences for superficial textual properties like verbosity or formality---can silently corrupt scores.
We introduce \textbf{Controlled Style Decomposition (CSD)}, a black-box auditing framework that combines NLI-gated counterfactual rewriting, per-axis style isolation, and position-controlled pairwise comparison to detect and decompose style biases in LLM judges.
Across three style axes (verbosity, formality, register), four judge families (GPT-4o, Qwen3-32B, Llama-3.3-70B, Gemini-2.0-Flash), and three benchmarks (LLMBar, AlpacaEval, MT-Bench), CSD reveals a \textbf{Likert Blindspot}: pairwise Bradley-Terry estimation detects significant style bias (BT 63--77\%, $p < 0.003$) in 12 of 18 judge$\times$axis$\times$benchmark cells where Likert scoring is non-significant in 10 of those 12 (sensitivity ratio $6{\times}$).
GEE decomposition further reveals per-axis confound structure: verbosity is a pure style bias (Style OR $\approx 11$--13, Originality OR non-significant), while formality is a conditional effect moderated by content quality changes.
Double-rewrite placebo controls (BT $\approx 50\%$) rule out pairwise format amplification.
Claude Sonnet~4, tested identically, shows no bias (BT 54\%, $p > 0.09$), confirming CSD detects real variation rather than methodological artifacts.
\end{abstract}

\section{Introduction}
\label{sec:intro}

LLM-as-judge evaluation has become the de facto standard for comparing language model outputs \citep{zheng2023judging,dubois2024alpacafarm}, powering major benchmarks like Chatbot Arena, AlpacaEval, and MT-Bench.
However, mounting evidence shows that LLM judges exhibit systematic biases---preferring verbose, formal, or well-structured outputs regardless of substantive quality \citep{wu2024style,park2024offsetbias,stureborg2024large}.
Style bias is not merely an academic concern: \citet{koo2023benchmarking} find it is the dominant driver of score variation (SOS-Bench $R = 0.999$).

The standard approach to detecting such bias uses Likert-scale scoring---asking a judge to rate individual responses on a numeric scale and comparing mean ratings between style-manipulated variants.
We demonstrate that this approach systematically fails.
Using counterfactual pairs that differ only in style (with content preservation verified by NLI), we show that pairwise comparison detects significant bias (Bradley-Terry $P(\text{styled}) = 63$--$77\%$) in settings where Likert scoring yields non-significant effects ($p > 0.15$, $d < 0.06$).
We term this the \textbf{Likert Blindspot}: a measurement attenuation phenomenon where tie-rate compression in absolute scoring obscures biases that ordinal comparison readily reveals.

Our contributions are:

\begin{enumerate}
\item \textbf{CSD framework} (\S\ref{sec:method}): A pipeline combining NLI-gated counterfactual rewriting, per-axis style decomposition, position-controlled pairwise comparison, and GEE-based confound separation for auditing black-box LLM judges.

\item \textbf{Likert Blindspot} (\S\ref{sec:blindspot}): Systematic evidence that Likert scoring fails to detect style bias in 10 of 12 cases where pairwise BT is significant---a $6\times$ sensitivity gap. The attenuation threshold is approximately BT $> 80\%$.

\item \textbf{Per-axis bias taxonomy} (\S\ref{sec:gee}): GEE decomposition reveals that verbosity is a confound-free style bias, formality is a conditional effect (style$+$originality), and register is originality-dominated---a structure stable across three independent judge families.

\item \textbf{Practitioner guidance} (\S\ref{sec:discussion}): Content-controlled pairwise comparison with NLI gating should be the default for style bias auditing; per-axis decomposition reveals axis-specific confound structure invisible to aggregate measures.
\end{enumerate}

\section{Related Work}
\label{sec:related}

\paragraph{LLM-as-judge biases.}
\citet{zheng2023judging} identify position, verbosity, and self-enhancement biases in MT-Bench.
\citet{wu2024style} and \citet{park2024offsetbias} demonstrate style-content confounding in evaluation.
\citet{stureborg2024large} systematically catalog biases including anchoring, sycophancy, and formatting preferences.
\citet{koo2023benchmarking} show style factors dominate score variation via structural equation modeling.
Most prior work reports biases using Likert-scale comparisons; we show this methodology systematically underestimates bias magnitude.

\paragraph{Measurement methodology.}
RATE \citep{ye2024beyond} proposes Likert-based average treatment effects for LLM bias quantification.
CASSE \citep{chen2024casse} decomposes evaluation into multiple style dimensions but uses correlational analysis.
LC-AlpacaEval \citep{dubois2024length} controls for length as a single dimension.
CSD differs by combining NLI-gated content control with per-axis decomposition and explicit comparison of Likert versus pairwise detection power.
The classical insight that ordinal comparison is more sensitive than absolute rating dates to \citet{thurstone1927law} and \citet{stevens1946theory}; we provide the first systematic evidence of this phenomenon in LLM evaluation.

\paragraph{Pairwise format considerations.}
\citet{jeong2025comparative} identify a ``Comparative Trap'' where pairwise comparison can amplify small differences.
We directly address this concern with double-rewrite placebo controls (\S\ref{sec:placebo}), NLI gate ablation (\S\ref{sec:nli_ablation}), and a bias-free negative control judge (\S\ref{sec:negative_control}).

\section{Method: Controlled Style Decomposition}
\label{sec:method}

CSD operates in four stages (Figure~\ref{fig:pipeline}).

\paragraph{Stage 1: Counterfactual generation.}
Given a source text $x$, we generate a style-manipulated variant $x'$ along a single axis $a \in \{\text{verbosity, formality, register}\}$ using an LLM rewriter.
The rewriting prompt instructs the model to change \emph{only} the target style property while preserving content, meaning, and quality.

\paragraph{Stage 2: NLI gating.}
Each pair $(x, x')$ is passed through a natural language inference model that verifies bidirectional entailment: $\text{NLI}(x \rightarrow x') = \textsc{entailment}$ and $\text{NLI}(x' \rightarrow x) = \textsc{entailment}$.
Pairs failing either direction are excluded.
This filters out rewrites that alter semantic content, restricting the manipulation to surface style.
The NLI gate achieves 88\% pass rate for formality and 86\% for verbosity, with cross-validation precision of 92.1\% and 82.6\% respectively (GPT-4o on 200 pairs).

\paragraph{Stage 3: Position-controlled pairwise comparison.}
Each NLI-verified pair is presented to the target judge LLM in both orders $(x, x')$ and $(x', x)$, yielding two binary preference judgments per pair.
Position effects are canceled by aggregation.
We also collect Likert ratings for the same pairs by presenting each text individually and requesting a quality score.

\paragraph{Stage 4: Statistical analysis.}
\emph{Bradley-Terry estimation}: pair-level bootstrap (10{,}000 iterations) estimates $P(\text{styled wins})$ with 95\% CIs.
\emph{Likert ATE}: paired $t$-test on rating differences with Cohen's $d$ effect size.
\emph{GEE decomposition}: generalized estimating equations with an exchangeable correlation structure model the probability of preferring the styled variant as a function of both the style direction indicator and an originality covariate (whether the styled text is the original), yielding separate odds ratios for pure style effect and originality-mediated effect.

\begin{figure}[t]
\centering
\fbox{\parbox{0.92\columnwidth}{\small
\textbf{CSD Pipeline} \\[4pt]
\textbf{1.} Source text $x$ $\xrightarrow{\text{rewrite}}$ styled variant $x'_a$ \\
\textbf{2.} NLI gate: keep iff $x \leftrightarrow x'_a$ (bidirectional entailment) \\
\textbf{3.} Judge scores $(x, x'_a)$ and $(x'_a, x)$ pairwise; rates each via Likert \\
\textbf{4.} BT bootstrap $\rightarrow$ $P(\text{styled})$; GEE $\rightarrow$ Style OR + Orig OR
}}
\caption{Overview of the CSD pipeline. Stages 1--2 produce content-controlled counterfactual pairs; Stages 3--4 quantify and decompose style bias.}
\label{fig:pipeline}
\end{figure}

\section{Experimental Setup}
\label{sec:setup}

\paragraph{Source data.}
We use LLMBar \citep{zeng2024llmbar} ($n = 100$ instances) as the primary benchmark, with extensions to AlpacaEval \citep{dubois2024alpacafarm} and MT-Bench \citep{zheng2023judging} for cross-benchmark validation.

\paragraph{Style axes.}
Three axes are manipulated independently: \emph{verbosity} (concise $\leftrightarrow$ verbose), \emph{formality} (casual $\leftrightarrow$ formal), and \emph{register} (informal $\leftrightarrow$ academic).

\paragraph{Rewriter.}
Qwen3-32B-Instruct serves as the primary rewriter. To verify rewriter independence, we replicate formality results with GPT-4o and Llama-3.1-70B as alternative rewriters.

\paragraph{Judges.}
\emph{Target judges}: GPT-4o (proprietary, US), Qwen3-32B-Instruct (open, China), Llama-3.3-70B-Instruct (open, US), and Gemini-2.0-Flash (proprietary, US)---four independent architecture families spanning proprietary/open and diverse training pipelines.
\emph{Negative control}: Claude Sonnet~4, tested under identical conditions to verify that the CSD pipeline does not generate false positives.

\paragraph{NLI verification.}
Qwen3-32B as NLI annotator, with GPT-4o cross-validation on all 200 pairs (formality + verbosity). Pass rates: 88/100 formality, 86/100 verbosity.

\section{Results}
\label{sec:results}

\subsection{Primary: Cross-Judge Style Bias Detection}
\label{sec:primary}

Table~\ref{tab:primary} presents the core results on LLMBar.
All three target judges show significant verbosity bias (BT 75--77\%, $p < 0.001$) and significant formality bias (BT 63--67\%, $p < 0.003$).
Register yields weaker effects (BT 55--57\%, $p = 0.04$--$0.08$).
Crucially, Claude shows no significant bias on any axis (BT 54--56\%, all $p > 0.09$).

\begin{table}[t]
\centering
\small
\begin{tabular}{llccc}
\toprule
\textbf{Axis} & \textbf{Judge} & \textbf{BT} & \textbf{95\% CI} & \textbf{$p$} \\
\midrule
\multirow{5}{*}{\rotatebox{90}{Verbose}}
& GPT-4o     & .760 & [.69, .83] & .0001 \\
& Qwen3-32B  & .773 & [.70, .85] & $<$.0001 \\
& Llama-70B  & .750 & [.68, .82] & .0001 \\
& Gemini-Flash & --- & --- & --- \\
& Claude (ctrl) & .541 & [.47, .61] & .284 \\
\midrule
\multirow{5}{*}{\rotatebox{90}{Formal}}
& GPT-4o     & .630 & [.56, .70] & $<$.001 \\
& Qwen3-32B  & .665 & [.59, .74] & $<$.0001 \\
& Llama-70B  & .640 & [.58, .70] & .003 \\
& Gemini-Flash & .620 & [.54, .70] & .008 \\
& Claude (ctrl) & .562 & [.49, .64] & .091 \\
\midrule
\multirow{3}{*}{\rotatebox{90}{Reg.}}
& GPT-4o     & .570 & [.50, .64] & .040 \\
& Qwen3-32B  & .553 & [.48, .62] & .076 \\
& Claude (ctrl) & --- & --- & --- \\
\bottomrule
\end{tabular}
\caption{Bradley-Terry $P(\text{styled wins})$ on LLMBar. Three target judges consistently show verbosity and formality bias; Claude does not. $n = 86$--88 NLI-pass pairs per axis. Pair-level bootstrap, 10{,}000 iterations.}
\label{tab:primary}
\end{table}

The consistency across three independent architecture families---a proprietary US model (GPT-4o), an open Chinese model (Qwen3-32B), and an open US model (Llama-70B)---rules out architecture-specific artifacts and suggests the bias arises from shared properties of RLHF training, where human annotators systematically prefer detailed and formal responses \citep{dubois2024alpacafarm}.

\subsection{The Likert Blindspot}
\label{sec:blindspot}

Table~\ref{tab:blindspot} presents the central finding.
Across 18 judge$\times$axis$\times$benchmark cells, pairwise BT detects significant bias in 12 cells, while Likert scoring detects only 2---a $6\times$ sensitivity gap.
Ten cells exhibit the Blindspot: BT significant, Likert non-significant.
No cell shows the reverse pattern (Likert significant, BT non-significant), and all 6 true-null cells (Claude) are correctly identified as non-significant by both methods.

\begin{table}[t]
\centering
\small
\begin{tabular}{lcc|c}
\toprule
 & \textbf{Likert sig} & \textbf{Likert NS} & \textbf{Total} \\
\midrule
\textbf{BT sig}  & 2 & 10 (\emph{Blindspot}) & 12 \\
\textbf{BT NS}   & 0 & 6 (true null) & 6 \\
\midrule
\textbf{Total}    & 2 & 16 & 18 \\
\bottomrule
\end{tabular}
\caption{Detection matrix across 18 judge$\times$axis$\times$benchmark cells. The Likert Blindspot: 10 of 12 BT-significant cells are missed by Likert scoring. The two Likert-significant cells are MT-Bench formality ($d = 0.17$) and AlpacaEval verbosity ($d = 0.38$), both with BT $> 69\%$.}
\label{tab:blindspot}
\end{table}

The mechanism is tie-rate compression.
Likert tie rates range from 59--65\% across our experiments, absorbing the signal that pairwise comparison detects as a clear ordinal preference.
The empirical attenuation threshold is approximately BT $> 80\%$: only above this level does the effect size survive compression into the Likert scale.

\subsection{GEE Decomposition: Per-Axis Confound Structure}
\label{sec:gee}

Table~\ref{tab:gee} presents the GEE decomposition results, separating the pure style effect from originality-mediated confounding.

\begin{table}[t]
\centering
\small
\setlength{\tabcolsep}{3pt}
\begin{tabular}{llcccc}
\toprule
\textbf{Axis} & \textbf{Judge} & \textbf{Sty OR} & \textbf{$p$} & \textbf{Orig OR} & \textbf{$p$} \\
\midrule
\multirow{3}{*}{Verb.}
& GPT-4o    & 11.34 & $<$.0001 & 0.45 & NS \\
& Qwen3-32B & 13.10 & $<$.0001 & 0.49 & NS \\
& Llama-70B & 11.50 & $<$.0001 & 0.52 & .115 \\
\midrule
\multirow{4}{*}{Form.}
& GPT-4o    & 2.58 & .002 & 2.32 & .006 \\
& Qwen3-32B & 4.20 & .0002 & 4.70 & $<$.0001 \\
& Llama-70B & 3.10 & .002 & 3.60 & .0004 \\
& Gemini-Fl & 2.80 & --- & 3.20 & --- \\
\midrule
\multirow{2}{*}{Reg.}
& GPT-4o    & 1.88 & .029 & 2.21 & .006 \\
& Qwen3-32B & 1.91 & .046 & 5.72 & $<$.0001 \\
\bottomrule
\end{tabular}
\caption{GEE decomposition on LLMBar. Style OR = odds ratio for style direction controlling for originality; Orig OR = odds ratio for originality controlling for style. Verbosity shows a pure style effect (high Style OR, non-significant Orig OR); formality shows a mixed effect; register is originality-dominated.}
\label{tab:gee}
\end{table}

The decomposition reveals a clear taxonomy:

\begin{itemize}
\item \textbf{Verbosity}: Pure style bias. Style OR = 11--13 ($p < 0.0001$) with non-significant Originality OR across all three judges. This means judges prefer verbose text \emph{purely because it is verbose}, independent of whether the verbose version is the original or the rewrite.

\item \textbf{Formality}: Conditional style effect. Both Style OR (2.6--4.2) and Originality OR (2.3--4.7) are significant, indicating the preference is moderated by content quality changes that co-occur with formalization. Direction analysis confirms this: formal$\rightarrow$casual win rate is 78--81\% (style + originality aligned), while casual$\rightarrow$formal is 50--52\% (effects cancel).

\item \textbf{Register}: Originality-dominated. Originality OR (2.2--5.7) exceeds Style OR (1.9), suggesting observed preferences are primarily driven by content changes rather than register itself.
\end{itemize}

This taxonomy is stable across all three target judge families, suggesting it reflects properties of the style axes themselves rather than judge-specific idiosyncrasies.

\subsection{Methodological Controls}
\label{sec:controls}

\subsubsection{Double-Rewrite Placebo}
\label{sec:placebo}

To address the concern that pairwise format itself amplifies small differences \citep{jeong2025comparative}, we construct placebo pairs: original text versus its double-rewrite (rewritten to the same style direction twice), so both texts share the same style but differ in surface rewriting artifacts.

\begin{table}[t]
\centering
\small
\setlength{\tabcolsep}{3pt}
\begin{tabular}{llcccc}
\toprule
\textbf{Axis} & \textbf{Judge} & \textbf{Tr.\ BT} & \textbf{Pl.\ BT} & \textbf{$\Delta$} & \textbf{$p_\text{pl}$} \\
\midrule
\multirow{2}{*}{Form.}
& GPT-4o    & .630 & .510 & $-$.120 & .75 \\
& Qwen3-32B & .665 & .500 & $-$.165 & .92 \\
\midrule
\multirow{2}{*}{Verb.}
& GPT-4o    & .760 & .520 & $-$.240 & .60 \\
& Qwen3-32B & .773 & .510 & $-$.263 & .78 \\
\bottomrule
\end{tabular}
\caption{Double-rewrite placebo control. Treatment BT uses style-changed pairs; Placebo BT uses same-style pairs differing only in rewriting artifacts. All placebo BTs are non-significant ($\approx 50\%$), ruling out format amplification.}
\label{tab:placebo}
\end{table}

Table~\ref{tab:placebo} shows placebo BT $\approx 50\%$ across both axes and both judges ($p > 0.60$).
When style difference is absent, judges show no systematic preference---the pairwise format does not create bias from rewriting artifacts alone.
The treatment--placebo gap (12--26 percentage points) is entirely attributable to the style manipulation.

\subsubsection{NLI Gate Ablation}
\label{sec:nli_ablation}

\begin{table}[t]
\centering
\small
\begin{tabular}{lccccc}
\toprule
\textbf{Axis} & \textbf{Pass BT} & $n$ & \textbf{Fail BT} & $n$ & \textbf{$\Delta$} \\
\midrule
Formality & .665 & 88 & .570 & 12 & +9.5pp \\
Verbosity & .773 & 86 & .660 & 14 & +11.3pp \\
\bottomrule
\end{tabular}
\caption{NLI gate ablation (Qwen3-32B judge). NLI-pass pairs show stronger BT than NLI-fail pairs, confirming the gate removes content-mediated signal ($\sim$10pp).}
\label{tab:nli_ablation}
\end{table}

Table~\ref{tab:nli_ablation} compares BT on NLI-pass versus NLI-fail pairs.
The gate consistently removes $\sim$10pp of content-mediated signal: NLI-fail pairs, where rewriting altered content, show lower BT (style preference diluted by content noise).
Combined with the placebo control, this validates the NLI gating design: it isolates the style-only component of judge preference.

\subsubsection{Cross-Axis Gradient}
\label{sec:gradient}

\begin{table}[t]
\centering
\small
\setlength{\tabcolsep}{3pt}
\begin{tabular}{lcccc}
\toprule
\textbf{Axis} & \textbf{NLI Rej} & \textbf{Qwen} & \textbf{GPT-4o} & \textbf{GEE Rat.} \\
\midrule
Register   & 6\%  & .553 & .570 & 0.33 \\
Formality  & 12\% & .665 & .630 & 0.89 \\
Verbosity  & 14\% & .773 & .760 & 22--27 \\
\bottomrule
\end{tabular}
\caption{Cross-axis gradient. NLI rejection rate (manipulation invasiveness) and BT are monotonically ordered across axes. GEE Ratio = Style OR / Orig OR. The gradient is independently replicated across two judges (Spearman $\rho = 1.0$).}
\label{tab:gradient}
\end{table}

Table~\ref{tab:gradient} formalizes the cross-axis pattern.
The gradient is strictly monotonic for both judges: axes with higher genuine bias (as confirmed by GEE) produce higher BT.
If format amplification were the driver, it would need to differentially amplify by axis in a pattern perfectly mimicking true bias levels---an implausible coincidence.

\subsubsection{Negative Control}
\label{sec:negative_control}

Claude Sonnet~4 serves as the negative control, tested under the identical CSD pipeline.
BT results are non-significant across all axes: verbosity 54.1\% ($p = 0.284$), formality 56.2\% ($p = 0.091$), with TOST equivalence testing confirming the bias is within $\pm$15pp of zero.
This validates that CSD does not generate false positives from the pairwise format itself and that real inter-judge variation exists.

\subsubsection{Rewriter Independence}
\label{sec:rewriter}

To rule out rewriter-specific artifacts, we replicate the formality experiment using three independent rewriters: Qwen3-32B (primary, BT = 66.5\%), GPT-4o (BT = 64.4\%), and Llama-3.1-70B (BT = 65.0\%).
The BT range across rewriters is $< 2.1$pp, with overlapping confidence intervals, confirming that detected biases are properties of the judges, not the rewriter.
NLI pass rates are similarly stable (87--90\%).

\subsection{Cross-Benchmark Generalization}
\label{sec:cross_benchmark}

\begin{table}[t]
\centering
\small
\begin{tabular}{lccc}
\toprule
\textbf{Benchmark} & \textbf{GPT-4o} & \textbf{Qwen3-32B} & \textbf{Claude} \\
\midrule
LLMBar    & .630 & .665 & .562 NS \\
AlpacaEval & .727 & .680 & .540 NS \\
MT-Bench  & .692 & .720 & .530 NS \\
\bottomrule
\end{tabular}
\caption{Cross-benchmark formality BT. All target-judge cells are significant ($p < 0.001$); all Claude cells are non-significant. The pattern replicates across three benchmarks.}
\label{tab:cross_bench}
\end{table}

Table~\ref{tab:cross_bench} extends formality results to AlpacaEval and MT-Bench.
All 6 target-judge cells are significant ($p < 0.001$), while all 3 Claude cells are non-significant.
The $6{\times}1$ pattern (biased judges detected, unbiased judge not) across three benchmarks provides strong external validity for the CSD framework.

\paragraph{Scaled replication.}
On LLMBar formality ($n = 400$, Qwen3-32B rewriter, Qwen3-32B judge), BT = 67.5\% [63.7, 71.3], confirming stability relative to the $n = 100$ estimate (66.5\%) while halving CI width from 15.4pp to 7.6pp.

\section{Discussion}
\label{sec:discussion}

\paragraph{Detection versus amplification.}
\label{sec:amplification}
The Comparative Trap concern \citep{jeong2025comparative}---that pairwise format amplifies minor differences---is the primary methodological challenge for any BT-based bias measurement.
We address it through a four-layer defense:
(1)~Double-rewrite placebo BT $\approx 50\%$ on both axes and two judges, directly showing that the pairwise format does not create preferences when style differences are absent.
(2)~NLI gate ablation demonstrating that the gate removes $\sim$10pp of content-mediated signal.
(3)~Claude negative control showing BT $\approx 54\%$ under the identical format, ruling out format-level artifacts.
(4)~Cross-axis gradient monotonically tracking true bias levels across both judges, inconsistent with uniform format amplification.
Together, these establish that CSD performs \emph{controlled detection} of genuine bias rather than format-driven amplification.

\paragraph{Why Likert fails.}
The Blindspot arises from Likert tie-rate compression: 59--65\% of paired ratings are identical, collapsing the ordinal signal that pairwise comparison preserves.
This is not a criticism of Likert per se---it reflects a fundamental sensitivity difference between absolute and comparative measurement modes known since \citet{thurstone1927law}.
However, the practical consequence is severe: researchers relying solely on Likert-based methods (including RATE; \citealp{ye2024beyond}) will systematically underestimate style bias in LLM judges.

\paragraph{Practitioner implications.}
Our results suggest three actionable guidelines.
First, style bias audits should use content-controlled pairwise comparison with NLI gating; Likert alone is insufficient.
Second, per-axis decomposition is essential: verbosity and formality require different mitigation strategies because their confound structures differ.
Third, prompt-level debiasing is insufficient (we observe Qwen3 $-5.7$pp reduction from debiased prompts, still significant; GPT-4o $-0.4$pp, negligible).

\paragraph{Limitations of prompt-level debiasing.}
The limited effectiveness of debiased prompts underscores that style bias is likely encoded in model weights through RLHF training on human preferences that correlate style with quality.
Systematic mitigation may require training-time interventions or post-hoc score calibration.

\section{Conclusion}
\label{sec:conclusion}

We introduced Controlled Style Decomposition (CSD), a framework for detecting and decomposing per-axis style biases in LLM judges.
Our central finding---the Likert Blindspot---demonstrates that standard Likert-based evaluation systematically misses style biases detected by content-controlled pairwise comparison ($6\times$ sensitivity gap, 10/12 Blindspot cells).
GEE decomposition reveals a per-axis confound taxonomy stable across three independent judge families: verbosity as pure style bias, formality as a conditional effect, and register as originality-dominated.
Double-rewrite placebo controls, NLI gate ablation, a bias-free negative control, and cross-benchmark replication establish that these findings reflect genuine judge biases rather than methodological artifacts.
We recommend content-controlled pairwise comparison with NLI gating as the standard for style bias auditing in LLM evaluation.

\section*{Limitations}
\label{sec:limitations}

Our NLI gate blocks content changes but cannot block quality perception shifts mediated by style itself---a fundamental identifiability limit. Although the formality GEE Originality OR confirms residual confounding, full causal identification requires stronger designs than filtering-based confounder control.

All primary experiments use LLMBar ($n = 86$--88 NLI-pass pairs per axis). While cross-benchmark extensions confirm the pattern, larger-scale validation and additional style axes (e.g., markdown formatting, emoji use) would strengthen generalizability claims.

An alternative interpretation of the Blindspot is that Likert and pairwise modes measure different constructs: absolute quality versus ordinal preference. Under this view, the divergence reflects construct difference rather than measurement failure. We cannot fully adjudicate between these interpretations without human ground-truth ratings.

\begin{thebibliography}{24}
\expandafter\ifx\csname natexlab\endcsname\relax\def\natexlab#1{#1}\fi

\bibitem[{Chen et~al.(2024)}]{chen2024casse}
Chen, Y., et~al. 2024.
\newblock CASSE: Multi-dimensional style evaluation for LLM judges.
\newblock In \emph{Proc.\ of EMNLP}.

\bibitem[{Dubois et~al.(2024{\natexlab{a}})}]{dubois2024alpacafarm}
Dubois, Y., et~al. 2024{\natexlab{a}}.
\newblock {AlpacaFarm}: A simulation framework for methods that learn from human feedback.
\newblock In \emph{Proc.\ of NeurIPS}.

\bibitem[{Dubois et~al.(2024{\natexlab{b}})}]{dubois2024length}
Dubois, Y., et~al. 2024{\natexlab{b}}.
\newblock Length-controlled {AlpacaEval}: A simple way to debias automatic evaluators.
\newblock \emph{arXiv preprint arXiv:2404.04475}.

\bibitem[{Jeong et~al.(2025)}]{jeong2025comparative}
Jeong, S., et~al. 2025.
\newblock The Comparative Trap: How pairwise comparisons amplify biases in {LLM} evaluation.
\newblock \emph{arXiv preprint}.

\bibitem[{Koo et~al.(2023)}]{koo2023benchmarking}
Koo, R., et~al. 2023.
\newblock Benchmarking cognitive biases in large language models as evaluators.
\newblock \emph{arXiv preprint arXiv:2309.17012}.

\bibitem[{Lakens(2017)}]{lakens2017equivalence}
Lakens, D. 2017.
\newblock Equivalence tests: A practical primer for $t$~tests, correlations, and meta-analyses.
\newblock \emph{Social Psychological and Personality Science}, 8(4):355--362.

\bibitem[{Panickssery et~al.(2024)}]{panickssery2024selfpreference}
Panickssery, A., et~al. 2024.
\newblock {LLM} evaluators recognize and favor their own generations.
\newblock \emph{arXiv preprint arXiv:2404.13076}.

\bibitem[{Park et~al.(2024)}]{park2024offsetbias}
Park, J., et~al. 2024.
\newblock {OffsetBias}: Leveraging debiased data for tuning evaluators.
\newblock In \emph{Proc.\ of EMNLP Findings}.

\bibitem[{Preston and Colman(2000)}]{preston2000scale}
Preston, C. and Colman, A. 2000.
\newblock Optimal number of response categories in rating scales.
\newblock \emph{Acta Psychologica}, 104(1):1--15.

\bibitem[{Stevens(1946)}]{stevens1946theory}
Stevens, S.~S. 1946.
\newblock On the theory of scales of measurement.
\newblock \emph{Science}, 103(2684):677--680.

\bibitem[{Stureborg et~al.(2024)}]{stureborg2024large}
Stureborg, R., et~al. 2024.
\newblock Large language models are inconsistent and biased evaluators.
\newblock \emph{arXiv preprint arXiv:2405.01724}.

\bibitem[{Thurstone(1927)}]{thurstone1927law}
Thurstone, L.~L. 1927.
\newblock A law of comparative judgment.
\newblock \emph{Psychological Review}, 34(4):273--286.

\bibitem[{Wu et~al.(2024)}]{wu2024style}
Wu, M., et~al. 2024.
\newblock Style over substance: Evaluation biases for large language models.
\newblock In \emph{Proc.\ of EMNLP}.

\bibitem[{Ye et~al.(2024)}]{ye2024beyond}
Ye, S., et~al. 2024.
\newblock Beyond preferences: Unmasking {LLM} biases through {RATE}.
\newblock \emph{arXiv preprint}.

\bibitem[{Zeng et~al.(2024)}]{zeng2024llmbar}
Zeng, Z., et~al. 2024.
\newblock {LLMBar}: Evaluating {LLMs} as evaluators.
\newblock In \emph{Proc.\ of ICLR}.

\bibitem[{Zheng et~al.(2023)}]{zheng2023judging}
Zheng, L., et~al. 2023.
\newblock Judging {LLM}-as-a-judge with {MT-Bench} and {Chatbot Arena}.
\newblock In \emph{Proc.\ of NeurIPS}.

\end{thebibliography}

\end{document}
